% !TEX program = platex
% 数学ノート集・共通テンプレート v2（積分論2026年度第1部の版面を参考）
% A5判。章内で定義・定理・例・注意等の番号を共通化。
% コンパイル： platex -interaction=nonstopmode math_note_template_v2.tex (2回)
%              dvipdfmx math_note_template_v2.dvi
\documentclass[a5paper,10pt,onecolumn]{jsarticle}
\usepackage[dvipdfmx,top=22mm,bottom=20mm,inner=19mm,outer=18mm,headheight=13pt,headsep=6mm]{geometry}
\usepackage{amsmath,amssymb,amsthm,mathtools}
\usepackage{enumitem,fancyhdr,titlesec}
\usepackage[dvipdfmx,bookmarks=true,colorlinks=false]{hyperref}
\usepackage{pxjahyper}

% ----- 差し替える文書情報 -----
\newcommand{\NoteTitle}{数値計算の誤差と安定性}
\newcommand{\NoteSeries}{数学ノート集}

% ----- 文字・組版 -----
% 参考PDFの質素な明朝＋Computer Modernの印象を引き継ぎながら，
% A5判で窮屈にならないよう本文の行間を少し確保する．
\linespread{1.13}
\setlength{\parindent}{1zw}
\setlength{\parskip}{0pt}
\setlength{\abovedisplayskip}{6pt plus 2pt minus 2pt}
\setlength{\belowdisplayskip}{6pt plus 2pt minus 2pt}
\setlength{\abovedisplayshortskip}{3pt plus 2pt minus 1pt}
\setlength{\belowdisplayshortskip}{4pt plus 2pt minus 1pt}
\setlist[enumerate]{topsep=3pt,itemsep=1pt,parsep=0pt,leftmargin=2.4zw}
\titleformat{\section}{\normalfont\normalsize\bfseries}{\thesection}{0.8zw}{}
\titlespacing*{\section}{0pt}{2.0ex plus .4ex}{.9ex}
\titleformat{\subsection}{\normalfont\normalsize\bfseries}{\thesubsection}{0.7zw}{}
\titlespacing*{\subsection}{0pt}{1.5ex plus .3ex}{.6ex}
\raggedbottom

% ----- 柱とページ番号：ページ上部、罫線なし -----
\fancypagestyle{mathnote}{%
  \fancyhf{}%
  \fancyhead[L]{\footnotesize \NoteSeries\ （\NoteTitle）}%
  \fancyhead[R]{\footnotesize \thepage}%
  \renewcommand{\headrulewidth}{0pt}%
  \renewcommand{\footrulewidth}{0pt}%
}
\pagestyle{mathnote}

% ----- 節内通し番号の定義・定理群 -----
\newtheoremstyle{mathnotetheorem}
 {4.5pt}{4.5pt}{\normalfont}{}%
 {\bfseries}{}{0.65zw}%
 {\thmname{#1}\thmnumber{ #2}\thmnote{（#3）}}
\theoremstyle{mathnotetheorem}
\newtheorem{theorem}{定理}[section]
\newtheorem{definition}[theorem]{定義}
\newtheorem{lemma}[theorem]{補題}
\newtheorem{proposition}[theorem]{命題}
\newtheorem{corollary}[theorem]{系}
\newtheorem{example}[theorem]{例}
\newtheorem{remark}[theorem]{注意}
\renewcommand{\proofname}{証明}
\renewcommand{\qedsymbol}{$\square$}
\makeatletter
\renewenvironment{proof}[1][\proofname]{%
  \par\pushQED{\qed}\normalfont\topsep4\p@\@plus2\p@\relax
  \trivlist\item[\hskip\labelsep\bfseries #1]\ignorespaces
}{\popQED\endtrivlist\@endpefalse}
\makeatother

% ----- 使うノートで適宜拡張する記号 -----
\newcommand{\R}{\mathbb{R}}
\newcommand{\norm}[1]{\lVert #1\rVert}
\newcommand{\abs}[1]{\lvert #1\rvert}

\begin{document}
\section*{到達目標と前提}
数値式の数学的な条件数と計算法の丸め誤差を区別し，
桁落ちが起こる式を自分で安定な形へ変形できるようになる．
前提は微分法と浮動小数点の四則演算の基礎である．

\section{数値計算に現れる三種類の誤差}
近似計算では，入力データの誤差，算法による丸め誤差，
数学的問題の摂動に対する感度を区別しなければならない．
\begin{definition}
真の値 $x\in\mathbb R$ と近似値 $\widehat x$ の
絶対誤差を $|\widehat x-x|$，
$x\ne0$ のとき相対誤差を
$|\widehat x-x|/|x|$ と定める．
\end{definition}
\begin{proposition}[相対誤差と零点]
$x$ が零に近いとき，小さい絶対誤差から
小さい相対誤差は一般には従わない．
\end{proposition}
\begin{proof}
$x=\varepsilon^2$，$\widehat x=\varepsilon^2+\varepsilon$
と置く．$\varepsilon>0$ を零へ近付けると
絶対誤差は $\varepsilon\to0$ である一方，
相対誤差は $1/\varepsilon\to\infty$ となる．
\end{proof}
\section{浮動小数点演算の数学的モデル}
以下では基数 $\beta\ge2$，有効桁数 $t\ge1$ の
正規化浮動小数点を考え，指数の範囲は十分広いとする．
オーバーフロー，アンダーフローは明記しない限り生じないと仮定する．
\begin{definition}
正規化した非零浮動小数点数は
\[
 \pm(d_0+d_1\beta^{-1}+\cdots+d_{t-1}\beta^{-(t-1)})
   \beta^e
\]
の形で表される．ここで $1\le d_0\le\beta-1$，
$0\le d_j\le\beta-1$，$e\in\mathbb Z$ とする．
実数 $x$ を最も近い浮動小数点数へ丸める操作を
$\operatorname{fl}(x)$ と書く．
\end{definition}
\begin{theorem}[基本丸めモデル]\label{thm:rounding}
丸める実数 $x$ が正規化範囲にあり，
最も近い数への正しい丸めが行われるとする．
このときある $\delta$ が存在して
\[
 \operatorname{fl}(x)=x(1+\delta),\qquad
 |\delta|\le u:=\frac12\beta^{1-t}
\]
が成立する．$x=0$ では $\delta=0$ と取れる．
\end{theorem}
\begin{proof}
$x>0$ としてよい．
$\beta^e\le x<\beta^{e+1}$ の区間で隣接する
浮動小数点数の間隔は $\beta^{e-t+1}$ である．
最も近い点への丸め誤差の絶対値は
その半分以下なので
\[
 |\operatorname{fl}(x)-x|
 \le\frac12\beta^{e-t+1}
 \le\frac12\beta^{1-t}x=ux.
\]
端点で隣の指数区間に丸める場合も隣接点間隔は
この上界以下であり同じ評価となる．
負の $x$ については対称性から従う．
\end{proof}
\begin{remark}
実際の計算機では指数の上限・下限，非正規化数，
NaN，無限大の取扱いがある．
定理\ref{thm:rounding}はそれらを除いた演算モデルであり，
全入力に対する無条件な保証ではない．
\end{remark}

\section{丸め誤差の蓄積}
本節では基本演算 $\circ\in\{+,-,\times,/\}$ の結果について
\[
 \operatorname{fl}(a\circ b)=(a\circ b)(1+\delta),
 \qquad |\delta|\le u
\]
が成立するものとし，零除算と指数範囲外を除外する．
\begin{lemma}[有限個の丸め因子]\label{lem:gamma}
$nu<1$ とし，$|\delta_j|\le u$（$1\le j\le n$）とする．
$\theta_n=\prod_{j=1}^n(1+\delta_j)-1$ と置けば
\[
 |\theta_n|\le\gamma_n:={nu\over1-nu}.
\]
\end{lemma}
\begin{proof}
全ての因子は正であるから
$(1-u)^n\le1+\theta_n\le(1+u)^n$．
Bernoulliの不等式により $(1-u)^n\ge1-nu$．
また $\log(1+u)\le u$ なので
$(1+u)^n\le e^{nu}$ である．
$0\le z<1$ では
$e^z\le(1-z)^{-1}$ が
$-\log(1-z)\ge z$ から従う．
従って
\[
 1-nu\le1+\theta_n\le(1-nu)^{-1}.
\]
負側の誤差は $nu$ 以下，
正側は $nu/(1-nu)$ 以下であり，
後者が前者以上なので結論を得る．
\end{proof}
\begin{theorem}[逐次加算の絶対誤差]\label{thm:sum}
実数 $x_1,\dots,x_n$（$n\ge2$）を
左から順に浮動小数点加算した値を $\widehat s$ とする．
各入力が正確に与えられ，途中の演算が基本丸めモデルに従い，
$(n-1)u<1$ なら
\[
 \left|\widehat s-\sum_{j=1}^nx_j\right|
 \le\gamma_{n-1}\sum_{j=1}^n|x_j|.
 \tag{1}
\]
\end{theorem}
\begin{proof}
最初の加算から順に丸め因子を
$(1+\delta_2),\dots,(1+\delta_n)$ とすると，
逐次加算の定義を展開して
\[
 \widehat s=
 x_1\prod_{k=2}^n(1+\delta_k)+
 \sum_{j=2}^n x_j\prod_{k=j}^n(1+\delta_k)
\]
を得る．
各係数の因子数は高々 $n-1$ である．
補題\ref{lem:gamma}を各積へ適用し，
三角不等式で誤差を抑えると式 (1) が従う．
\end{proof}
\begin{corollary}[加算の相対誤差]
$s=\sum_jx_j\ne0$ なら
\[
 {|\widehat s-s|\over|s|}
 \le\gamma_{n-1}{\sum_j|x_j|\over|\sum_jx_j|}.
\]
\end{corollary}
\begin{proof}
式 (1) の両辺を $|s|$ で割る．
\end{proof}
\begin{remark}
正負の大きい数が相殺すると
$\sum_j|x_j|/|\sum_jx_j|$ は非常に大きくなる．
これは丸め誤差の大きさだけでなく，
和という数学的問題の摂動感度が大きいことを示す．
\end{remark}
\begin{theorem}[二分木による和]\label{thm:pairwise}
$n=2^m$ 個の入力を二つずつ加え，
得た値をさらに二つずつ加える操作を $m$ 段階繰り返す．
その結果を $\widehat s_{\mathrm{pair}}$ とする．
$mu<1$ なら
\[
 \left|\widehat s_{\mathrm{pair}}-\sum_{j=1}^nx_j\right|
 \le\gamma_m\sum_{j=1}^n|x_j|.
\]
\end{theorem}
\begin{proof}
各入力 $x_j$ はただ一つの枝を通って最終和へ届き，
その枝では各段階ごとにちょうど一回の浮動小数点加算を受ける．
再帰的に展開した最終和では，
$x_j$ の係数は $m$ 個の丸め因子の積となる．
補題\ref{lem:gamma}をその積へ適用し，
全ての項の誤差を加えればよい．
\end{proof}
\begin{proposition}[内積の誤差]
入力の積を一回ずつ丸め，それらを左から順に加算した
$\widehat d$ を考える．$n\ge2$ かつ $nu<1$ なら
\[
 \left|\widehat d-\sum_{j=1}^n x_jy_j\right|
 \le\gamma_n\sum_{j=1}^n|x_jy_j|.
\]
\end{proposition}
\begin{proof}
各 $x_jy_j$ には積の丸め因子が一つ付き，
その後に高々 $n-1$ 個の加算の丸め因子が付く．
従って各項の係数は高々 $n$ 個の因子
$(1+\delta)$ の積である．
補題\ref{lem:gamma}と三角不等式から結論を得る．
\end{proof}

\section{問題の条件数}
\begin{definition}
開集合 $U\subset\mathbb R^m$ の写像
$F:U\to\mathbb R^n$ に対し，$x\in U$ で
$F$ が微分可能かつ $F(x)\ne0$ なら，
Euclidノルムについての相対条件数を
\[
 \kappa_F(x)={\|DF(x)\|\,\|x\|\over\|F(x)\|}
\]
と定める．
\end{definition}
\begin{proposition}[条件数の一次評価]
$x\ne0$，$F(x)\ne0$ とする．$h\to0$ のとき
\[
 {\|F(x+h)-F(x)\|\over\|F(x)\|}
 \le\kappa_F(x){\|h\|\over\|x\|}
 +o(\|h\|)
\]
が成立する．
\end{proposition}
\begin{proof}
微分可能性から
$F(x+h)-F(x)=DF(x)h+r(h)$，
$\|r(h)\|/\|h\|\to0$ である．
従って
\[
 \|F(x+h)-F(x)\|
 \le\|DF(x)\|\|h\|+o(\|h\|).
\]
両辺を $\|F(x)\|$ で割って定義を用いればよい．
\end{proof}
\begin{definition}
可逆な行列 $A\in\mathbb R^{n\times n}$ と
作用素ノルム $\|\cdot\|$ について
$\kappa(A)=\|A\|\|A^{-1}\|$ と定める．
\end{definition}
\begin{lemma}[逆行列の摂動]\label{lem:inverse}
$A$ が可逆で $\|A^{-1}E\|<1$ なら
$A+E$ は可逆であり，
\[
 \|(A+E)^{-1}\|
 \le{\|A^{-1}\|\over1-\|A^{-1}E\|}.
\]
\end{lemma}
\begin{proof}
$A+E=A(I+A^{-1}E)$ と書く．
$\|A^{-1}E\|<1$ なのでNeumann級数
$\sum_{k=0}^{\infty}(-A^{-1}E)^k$
は作用素ノルムで絶対収束する．
部分和へ $(I+A^{-1}E)$ を掛けると
$I-(-A^{-1}E)^{N+1}$ となり極限で $I$ を得るため，
これは $(I+A^{-1}E)^{-1}$ である．
ノルムを取ると幾何級数から
$\|(I+A^{-1}E)^{-1}\|\le
 1/(1-\|A^{-1}E\|)$ となる．
$(A+E)^{-1}=(I+A^{-1}E)^{-1}A^{-1}$ を使えばよい．
\end{proof}
\begin{theorem}[線形方程式の相対摂動評価]\label{thm:perturb}
$Ax=b$ とし，$A$ は可逆，$b\ne0$ とする．
$(A+E)(x+\Delta x)=b+\Delta b$ で，
\[
 \varepsilon_A={\|E\|\over\|A\|},\qquad
 \varepsilon_b={\|\Delta b\|\over\|b\|},\qquad
 \kappa(A)\varepsilon_A<1
\]
を仮定する．このとき
\[
 {\|\Delta x\|\over\|x\|}
 \le{\kappa(A)\over1-\kappa(A)\varepsilon_A}
 (\varepsilon_A+\varepsilon_b).
\]
\end{theorem}
\begin{proof}
差を取ると
$(A+E)\Delta x=\Delta b-Ex$ である．
補題\ref{lem:inverse}から
\[
 \|\Delta x\|
 \le{\|A^{-1}\|\over1-\|A^{-1}\|\|E\|}
 (\|\Delta b\|+\|E\|\|x\|).
\]
$Ax=b$ だから
$\|b\|\le\|A\|\|x\|$．
従って
$\|\Delta b\|/\|x\|\le
\varepsilon_b\|A\|$．
両辺を $\|x\|$ で割り，
$\|A^{-1}\|\|E\|=\kappa(A)\varepsilon_A$
を上から評価して結論を得る．
\end{proof}
\begin{example}[悪条件の線形方程式]
$A_\varepsilon=\operatorname{diag}(1,\varepsilon)$，
$0<\varepsilon\le1$ とする．
Euclid作用素ノルムについて
$\kappa(A_\varepsilon)=1/\varepsilon$ である．
\end{example}
\begin{proof}
$\|A_\varepsilon\|=1$，
$\|A_\varepsilon^{-1}\|=1/\varepsilon$ だからである．
$b=(1,\varepsilon)^T$ の解は $x=(1,1)^T$ だが，
右辺を $(1,\varepsilon+\delta)^T$ に変えると
解は $(1,1+\delta/\varepsilon)^T$ となる．
従って右辺の小さな誤差が
$1/\varepsilon$ 倍に増幅され得る．
\end{proof}

\section{前進安定性と後退安定性}
\begin{definition}
数学的問題を $F(x)$ とし，算法の出力を $\widehat y$ とする．
$\|\widehat y-F(x)\|$ を前進誤差という．
$\widehat y=F(x+\Delta x)$ となる摂動 $\Delta x$ の大きさを
後退誤差といい，そのような摂動を小さく取れる算法を
後退安定と呼ぶ．ここで「小さい」は問題の尺度と
単位丸め誤差 $u$ に相対的に定量化しなければならない．
\end{definition}
\begin{proposition}[後退誤差から前進誤差へ]
$F$ が $x$ で微分可能，$x,F(x)\ne0$ とする．
算法の出力が $\widehat y=F(x+\Delta x)$ を満たし
$\Delta x\to0$ なら
\[
 {\|\widehat y-F(x)\|\over\|F(x)\|}
 \le\kappa_F(x){\|\Delta x\|\over\|x\|}
 +o(\|\Delta x\|).
\]
\end{proposition}
\begin{proof}
条件数の一次評価を $\Delta x$ に適用する．
\end{proof}
\begin{remark}
後退安定な算法であっても条件数が大きい問題では
前進誤差は大きくなり得る．
例えば近い二数の減算は，丸めた差自体が非常に正確でも，
入力をわずかに変えた際の差の相対変化は大きくなり得る．
\end{remark}
\section{Horner法の誤差解析}
\begin{definition}
多項式 $p(x)=a_0+a_1x+\cdots+a_nx^n$ に対し
\[
 b_n=a_n,\qquad b_k=a_k+xb_{k+1}
 \quad(k=n-1,\ldots,0)
\]
によって $p(x)=b_0$ を求める算法をHorner法という．
\end{definition}
\begin{lemma}[Horner法の正しさ]
正確な実数演算でHorner法を行えば
\[
 b_k=\sum_{j=k}^n a_jx^{j-k}
\]
となり，特に $b_0=p(x)$ である．
\end{lemma}
\begin{proof}
$k=n$ では $b_n=a_n$ なので式は成立する．
$k+1$ で式が成り立つとすると
\[
 b_k=a_k+x\sum_{j=k+1}^n a_jx^{j-k-1}
 =\sum_{j=k}^na_jx^{j-k}.
\]
下向き帰納法により全ての $k$ で成り立つ．
\end{proof}
\begin{theorem}[Horner法の係数ごとの後退誤差]\label{thm:horner}
$x$ と全ての係数 $a_k$ が正確に与えられ，
Horner法の乗算と加算を別々に丸めるとする．
$2nu<1$，オーバーフロー・アンダーフローなしと仮定すると，
計算結果 $\widehat p$ に対し
\[
 \widehat p(x)=\sum_{k=0}^na_k(1+\theta_k)x^k,
 \qquad |\theta_k|\le\gamma_{2n}
\]
となる $\theta_k$ が存在する．従って
\[
 |\widehat p(x)-p(x)|
 \le\gamma_{2n}\sum_{k=0}^n|a_k||x|^k.
\]
\end{theorem}
\begin{proof}
実際の計算では
\[
 \widehat b_k=
 \bigl(a_k+x\widehat b_{k+1}(1+\delta_{k,1})\bigr)
 (1+\delta_{k,2}),\qquad
 |\delta_{k,j}|\le u
\]
であり，$\widehat b_n=a_n$ とする．
この再帰式を展開すれば，$a_kx^k$ の係数は
各段階の乗算・加算から生じる $1+\delta$ の積となる．
最上位係数から定数項に至るまで高々 $n$ 回の乗算と
$n$ 回の加算があるので，各係数に付く因子は高々 $2n$ 個である．
補題\ref{lem:gamma}よりその積は
$1+\theta_k$，$|\theta_k|\le\gamma_{2n}$ と書ける．
和全体から正確な多項式値を引いて三角不等式を適用すると
第二式を得る．
\end{proof}
\begin{remark}
Horner法の結論は，値 $\widehat p(x)$ が
各係数 $a_k$ を高々 $\gamma_{2n}$ の相対量だけ摂動させた
別の多項式の正確な値であることを示す．
ただし $p(x)$ が零に近い場合，値の相対前進誤差の
上界は大きくなることがある．
\end{remark}

\section{桁落ちを診断して安定な式へ変える}
\begin{example}[同じ関数の二つの計算法]
$x>0$ に対し
\[
 f(x)=\sqrt{1+x}-1
\]
を計算する．この式をそのまま計算すると，
$x$ が零に近いとき，近接する二数の差を求めるため
有効桁が失われる可能性がある．
分母を有理化すれば厳密に
\[
 f(x)=\frac{x}{\sqrt{1+x}+1}
\]
である．後者では近い二数の引算を避けられる．
なぜ元の問題自体が悪条件なのではないかを確認する．
$f'(x)=1/(2\sqrt{1+x})$ だから，
正の入力・出力に関する相対条件数は
\[
 \kappa_f(x)=\left|\frac{xf'(x)}{f(x)}\right|
 =\frac{\sqrt{1+x}+1}{2\sqrt{1+x}}.
\]
従って $1/2<\kappa_f(x)<1$（$x>0$）であり，
$x\downarrow0$ でも条件数は1に収束する．
問題が良条件でも計算式が不安定になり得る例である．
\end{example}
\begin{proposition}[引算による誤差増幅]
$x>0$ とし，正確な値 $\sqrt{1+x}$ に対して，
平方根を計算した近似値が
$\widehat s=\sqrt{1+x}+\varepsilon$ だったとする．
引算が正確に行われると仮定すると，
直接式で計算した値 $\widehat f=\widehat s-1$ の
相対誤差は厳密に
\[
 \frac{|\widehat f-f(x)|}{f(x)}
 =\frac{|\varepsilon|(\sqrt{1+x}+1)}{x}
\]
となる．
\end{proposition}
\begin{proof}
差 $\widehat f-f(x)$ は $\varepsilon$ に等しい．
一方，前例の有理化公式より
$f(x)=x/(\sqrt{1+x}+1)>0$ なので
両辺を $f(x)$ で割れば表示式が従う．
\end{proof}
\begin{remark}
平方根が丸められた後の引算にさらに丸めが発生すれば，
誤差は追加され得る．ここでは桁落ちの原因を切り出すため，
引算自体は正確であるとした．
有理化した式を浮動小数点で評価する場合にも，
除算や平方根の丸めの仮定を別途確認する必要がある．
\end{remark}
\section*{到達確認問題と解答}
\noindent\textbf{問題．}
$x>0$ のとき
$g(x)=1/(\sqrt{1+x}+1)$ を用いて
$f(x)=xg(x)$ と計算すると，
$x\downarrow0$ で $g(x)$ は何に収束するか．
また $x>0$ に対して
$0<f(x)<x/2$ を証明せよ．

\medskip\noindent\textbf{解答．}
$g(x)\to1/2$ である．
分母は $2$ より大きいから
$0<f(x)=x/(\sqrt{1+x}+1)<x/2$．
この評価から $f(x)$ が $x/2$ と同程度の
小さな値になることも分かる．

\end{document}
